\iffalse
\title{assignment}
\author{EE24BTECH11020}
\section{mcq-single}
\fi


%\begin{enumerate}

\item Let $f : \mathbb{R} \to \mathbb{R}$ be a continuous function. Then,

\begin{align*}
\lim_{x \to \frac{\pi}{4}} \frac{\frac{\pi}{4} \int_{2}^{\sec^2 x} f\brak{x} \, dx}{x^2 - \frac{\pi^2}{16}} \text{  is equal to:}
\end{align*}
\hfill (September 2021)


\begin{multicols}{4}
\begin{enumerate}
    \item $f\brak{2}$
    \item $2f\brak{2}$
    \item $2f\brak{\sqrt{2}}$
    \item $4f\brak{2}$
\end{enumerate}
\end{multicols}

\item $\cos^{-1}\brak{\cos\brak{-5}} + \sin^{-1}\brak{\sin\brak{6}} - \tan^{-1}\brak{\tan\brak{12}}$ is equal to:\hfill (September 2021)
\begin{multicols}{2}
    \begin{enumerate}
    \item $3\pi - 11$
    \item $4\pi - 9$
    \item $4\pi - 11$
    \item $3\pi + 1$
\end{enumerate}
\end{multicols}


\item Consider the system of linear equations:
\begin{align*}
-x + y + 2z &= 0 \\
3x - ay + 5z &= 1 \\
2x - 2y - az &= 7
\end{align*}
Let $S_1$ be the set of all $a \in \mathbb{R}$ for which the system is inconsistent, and $S_2$ the set for infinitely many solutions. Then:\hfill (September 2021)

\begin{multicols}{2}
    \begin{enumerate}
    \item $n\brak{S_1} = 2$, $n\brak{S_2} = 2$
    \item $n\brak{S_1} = 1$, $n\brak{S_2} = 0$
   \item $n\brak{S_1} = 2$, $n\brak{S_2} = 0$
    \item $n\brak{S_1} = 0$, $n\brak{S_2} = 2$
    \end{enumerate}
\end{multicols}

\item Let the acute angle bisector of the two planes
\begin{align*}
x - 2y - 2z + 1 &= 0 \\
2x - 3y - 6z + 1 &= 0
\end{align*}
be the plane $P$. Then which of the following points lies on $P$ ?\hfill (September 2021)
\begin{multicols}{2}
\begin{enumerate}
    \item $\brak{ 3, 1, -\frac{1}{2} }$ \\
    \item $\brak{0, 2, -4}$
     \item $\brak{ -2, 0, -\frac{1}{2} }$ \\
    \item $\brak{4, 0, -2}$
\end{enumerate}
\end{multicols}

\item Which of the following is equivalent to the Boolean expression $p$  $\land$ $\sim$$q$ ?\hfill (September 2021)
\begin{multicols}{4}
\begin{enumerate}
    \item $\sim$\brak{q \to p} 
    \item  $\sim\!p \to  \sim\!q$
    \item $\sim\brak{p \to  \sim\!q}$
    \item $\sim\!\brak{p \to q}$ 
\end{enumerate}
\end{multicols}

\item Two squares are chosen at random on a chessboard (see figure). The probability that they have a side in common is:\hfill (September 2021)
\begin{align}
\begin{tikzpicture}[scale=0.5]
    \draw[thin] (0, 0) grid (8, 8);
\end{tikzpicture}
\end{align}


\begin{multicols}{2}

\begin{enumerate}

    \item $\frac{2}{7}$ \\
    \item $\frac{1}{18}$ 
    \item $\frac{1}{7}$ \\
    \item $\frac{1}{9}$
\end{enumerate}
\end{multicols}

\item If $y = y\brak{x}$ is the solution curve of the differential equation
\begin{align*}
x^{2} \frac{dy}{dx} + y - \frac{1}{x} &= 0
\end{align*}
with $x>0$ and $y\brak{1}= 1$, then $y\brak{\frac{1}{2}}$ is equal to:\hfill (September 2021)
\begin{multicols}{2}
\begin{enumerate}
    \item $\frac{3}{2}-\frac{1}{\sqrt{e}}$\\
  
    \item $3 + e$
     \item $3 + \frac{1}{\sqrt{e}}$ \\
    \item $3 - e$
\end{enumerate}
\end{multicols}
\item If n is the number of solutions of the equation.\\
\begin{align*}
2 \cos{x}\brak{4\sin{\brak{\frac{\pi}{4} + x}}\sin{\brak{\frac{\pi}{4}-x}}-1}&=1,\\
\end{align*} $x \in \sbrak{0, \pi}$ and $S$ sum of all these solutions, Then ordered pair $(n, S)$ is:\hfill (September 2021)
\begin{multicols}{4}
\begin{enumerate}

    \item $(3, \frac{13\pi}{9})$
    \item $(2, \frac{2\pi}{3})$
    \item $(2, \frac{8\pi}{9})$
    \item $(3, \frac{5\pi}{3})$
\end{enumerate}
\end{multicols}
\item The function $f(x) = x^3 - 6x^2 + ax + b$ satisfies $f(2) = f(4) = 0$. Consider two statements:\\[10pt]
$\brak{S_1}$ There exist $x_1, x_2 \in \brak{2, 4}$ with $x_1 < x_2$, such that $f'\brak{x_1} = -1$ and $f'\brak{x_2} = 0$.\\[10pt]
$\brak{S_2}$ There exist $x_3, x_4 \in \brak{2, 4}$ with $x_3 < x_4$, such that $f$ is decreasing in $(2, x_4)$, increasing in $(x_4, 4)$, and $2f'(x_3) = \sqrt{3}f'(x_4)$.
Which of the following is correct?\hfill (September 2021)
\begin{multicols}{2}
\begin{enumerate}
    \item Both $\brak{S_1}$ and $\brak{S_2}$  are true
    \item $\brak{S_1}$ is false and $\brak{S_2}$  is true
    \item Both $\brak{S_1}$ and $\brak{S_2}$ are false
    \item $\brak{S_1}$ is true and $\brak{S_2}$ is false
\end{enumerate}
\end{multicols}

\item \begin{align*}
\text{ Let } J_{n,m} =  \int_0^\frac{1}{2} \frac{x^n}{x^m - 1} \, dx, \quad \forall n > m \text{ and } n, m \in \mathbb{N}.
\end{align*}
\text{Consider a matrix } $A = [a_{ij}]_{3 \times 3}$ \text{ where }
$a_{ij} = 
\begin{cases}
   J_{6+i, 3} - J_{i+3, 3} & i \leq j,\\
    0 & i > j.
    \end{cases}$
$\text{ Then } \abs{\text{adj}A^{-1}}  \text{ is:}$\hfill (September 2021)
    
\begin{multicols}{2}
\begin{enumerate}
    \item $(15)^2 \times 242$\\
    \item $(15)^2 \times 234$
    \item $(105)^2 \times 238$\\
    \item $(105)^2 \times 236$
\end{enumerate}
\end{multicols}

\item The area enclosed by the curves $y = \sin x - \cos x$ and $y = \abs{\cos x - \sin x}$ between $x = 0$ and $x = \frac{\pi}{2}$ is:\hfill (September 2021)
\begin{multicols}{2}
\begin{enumerate}
    \item $2\sqrt{2}(\sqrt{2}-1)$\\
    \item $2\sqrt{2}(1+\sqrt{2})$
    \item $4\sqrt{2}(\sqrt{2}-1)$\\
    \item $2\sqrt{2}(2+\sqrt{2})$
\end{enumerate}
\end{multicols}
\item The distance from the line $3y - 2z - 1 = 0 = 3x - z + 4$ to the point $(2, -1, 6)$ is:\hfill (September 2021)
\begin{multicols}{2}
\begin{enumerate}
    \item $2\sqrt{6}$\\
    \item $2\sqrt{5}$
    \item $4\sqrt{2}$\\
    \item $\sqrt{26}$
\end{enumerate}
\end{multicols}

\item Consider the parabola with vertex $\brak{\frac{1}{2}, \frac{3}{4}}$ and directrix $y = \frac{1}{2}$. Let $P$ be the point where the parabola\\ meets the line $x = -\frac{1}{2}$. If the normal to the parabola at $P$ intersects it again at the point $Q$, then\\ $(PQ)^2$ is equal to :\hfill (September 2021)
\begin{multicols}{2}
    \begin{enumerate}
    \item $\dfrac{75}{8}$
    \item $\dfrac{125}{16}$
    \item $\dfrac{25}{2}$
    \item $\dfrac{15}{2}$
\end{enumerate}
\end{multicols}


\item The number of pairs $\brak{a, b}$ of real numbers such that whenever $\alpha$ is a root of the equation $x^2 + ax + b = 0$, $\alpha^2 - 2$ is also a root is:\hfill (September 2021)
\begin{multicols}{4}
\begin{enumerate}
    \item 6
    \item 2
    \item 4
    \item 8
\end{enumerate}
\end{multicols}

\item Let $S_n = 1 \cdot (n-1) + 2 \cdot (n-2) + 3 \cdot (n-3) + \dots + (n-1) \cdot 1$, for $n \geq 4$.
\begin{align*} \text{The sum }
\sum_{n=4}^{\infty} \brak{\frac{2S_n}{n!} - \frac{1}{\brak{n-2}!}} \text{is equal to:}
\end{align*}\hfill (September 2021)
\begin{multicols}{2}
\begin{enumerate}
    \item $\frac{e-1}{3}$\\
    \item $\frac{e-2}{6}$
    \item $\frac{e}{3}$\\
    \item $\frac{e}{6}$
\end{enumerate}
\end{multicols}
